\noindent\makebox[0pt][l]{364}\hfill{\small\textbf{CORRESPONDENCE WITH BOHNENBERGER.}}\hfill\null\par\medskip

\begin{center}
[3].\\
\textsc{Gauss} to \textsc{Bohnenberger}.\quad Göttingen, 16~November 1823.
\end{center}

\par\noindent\hspace*{\parindent}I was pleased to learn something of your measurements from your letter, and nothing would please me more than if this reply were to prompt you to send me a somewhat fuller account of them.

\par\noindent\hspace*{\parindent}To begin with the heliotropes, I have had two quite different kinds made, based on quite different principles, but I find the second design the most advantageous.\ .\ .\ .\ .\ .\quad The small outlay [for the heliotropes] is saved many times over because, wherever the terrain otherwise permits, one can use the largest triangles and start work everywhere at once without first having to erect signals, which cost so much time and money. This is the least of the advantages, however, since the measurements can thereby attain a precision one can seldom count on with signals and church towers.
My worst triangles (relatively speaking) are those in which towers were the targets.
So much for the heliotropes.

\par\noindent\hspace*{\parindent}As for the measurements themselves, there is nothing I want more than to establish the connection with yours soon.
I have now connected my triangles with those of Electoral Hesse, and the enclosed drawing will give you an idea of the connection~[\textsuperscript{*})].
Only part of the Hessian triangles have actually been measured so far, but enough to connect all the points as far as Feldberg with mine, at least provisionally.
You know my triangles farther north from \textsc{Schumacher}'s Astronomischen Nachrichten, I, no.~24. The northern points added this year will be of no interest to you, at least for the moment.

\par\noindent\hspace*{\parindent}I now suppose that by way of this connection and of the Darmstadt and Bavarian triangles my measurements are linked to yours, but I have \textls[120]{none at all} of the latter yet and only some fragmentary data for the former.
If you have fuller data, you would greatly oblige me by communicating them, as well as the data for your own triangles.
I would only ask to receive the angles reduced to the centres but without adjustment for errors~\textsuperscript{**}).
I would like to compute everything
\par\smallskip\noindent\rule{.24\textwidth}{.4pt}\par\smallskip
{\footnotesize
\noindent\hspace*{1.5em}[\textsuperscript{*}) The copy of the original in the Gauss Archive is unaccompanied by a drawing.]\par
\noindent\hspace*{1.5em}\textsuperscript{**}) If need be, without any further results of the computations.\par
}
\clearpage
\noindent\null\hfill{\small\textbf{ON THE HANOVER TRIANGULATION.}}\hfill\makebox[0pt][r]{365}\par\medskip
\noindent by a uniform method.
I am all the more eager for the connection with Tübingen, Mannheim, and München because, \textls[120]{between ourselves}, the first measurements by \textsc{Schumacher} in Altona, where he has set up a fine little observatory with a \textsc{Reichenbach} meridian circle exactly like mine and those of \textsc{Soldner} and \textsc{Bessel}, gave an altitude of the pole differing by \(5^{\prime\prime}\) from the one transferred from Göttingen through the triangles.
In the computation I used the dimensions of the Earth that \textsc{Walbeck} derived from the full set of good arc measurements, with a flattening of \(\frac{1}{302{,}78}\).
Otherwise my methods of computation differ so completely from the ones in use that I cannot give you an idea of them in a letter.
Once my measurements are complete, I intend to develop these methods into a larger work and explain them by applying them to the Hanover measurements and those connected with them.
Our government is inclined to extend the Hanover measurements farther west, which would connect them with \textsc{Krayenhoff}'s and through those with the French and English ones.
\textsc{Krayenhoff}'s measurements, as you know, have been printed. The same should be done with all good triangulations. The large triangles belong, in a sense, to the entire educated world, present and future, all the more as they gradually become more closely interconnected.
\textsc{Müffling}'s triangles are already linked to the French arc measurement through \textsc{Tranchot}'s triangles. The latter, however, have not been printed and so are as good as nonexistent for the public, and General \textsc{von Müffling} himself has them only in an inadequate form (presumably the fatal chord angles already adjusted to \(180^{0}\)).
Unfortunately, I find people so disinclined to share information. I have gone to great lengths through official diplomatic and unofficial private channels to obtain from Paris the measurements \textsc{Epailly} made in Hanover in 1804 and 1805, but have received nothing but evasions, a bare list of station names, and a drawing of the triangles, apart from 2 or 3 numerical data that my measurements show to be decidedly grossly incorrect.
Let us make an exception.
I repeat my request that you send me your measurements and that you kindly share what you have of others' measurements connected with ours, insofar as I cannot obtain it directly, and I am happy to offer to reciprocate.
\par\clearpage
\noindent\makebox[0pt][l]{366}\hfill{\small\textbf{CORRESPONDENCE WITH BOHNENBERGER.}}\hfill\null\par\medskip
\noindent How splendid it would be if all the measurements extending across Europe, from Scotland to the Banat and from Copenhagen to Genua and Formentera, could one day be brought into a single connected system.
I would gladly do what I can to prepare the way for it, but when one is past the midpoint of one's life, the sooner one starts on so extensive a subject the better.

\par\noindent\hspace*{\parindent}\textsc{Gerling}'s measurements were made with a 12-inch \textsc{Ertel} theodolite exactly like mine and \textsc{Schumacher}'s\textsuperscript{*}).
In my measurements I have found that what I call the \textls[120]{mean} error in my paper in the latest Göttingen Commentationes, »Theoria combinationis observationum etc.«, is about \(=\frac{3\mathord{\overset{\prime\prime}{,}}5}{\surd n}\), \(n=\) the number of repetitions, when good and \textls[120]{less good} measurements from several stations are reckoned together.
With a very firm setup, very favourable (i.e.\ unshimmering) air, and exclusively heliotropic targets, however, it is considerably smaller.
All my measurements to date give 76 principal directions (38 forward and 38 in reverse), and the adjustment for errors showed the \textls[120]{mean} error of a principal direction to be \(=0\mathord{\overset{\prime\prime}{,}}47\).
They form a total of 26 triangles in which I measured all the angles, including several that form crossed quadrilaterals and pentagons. The adjustment was made without arbitrary decisions, without selection, and without exclusions, on strict principles of the calculus of probabilities, so that in the end everything fits together exactly.
Such crossed quadrilaterals would provide an excellent test for some surveys in which the sums of the angles are found to agree excellently everywhere, so that a triangle seldom fails to close by much over \(1^{\prime\prime}\). Yet that test (which is less within everyone's reach than computing a spherical excess), when it can be applied, sometimes shows quite conclusively that individual angles contain errors of \(2^{\prime\prime}\), \(3^{\prime\prime}\), or more.
The largest error in the sum of the unadjusted angles in my 26 triangles was \(2\mathord{\overset{\prime\prime}{,}}2\), where one direction went to the Michaelisthurm in Hamburg, which was very hard to sight and not entirely free from phase effects in sunshine. The next largest was \(1\mathord{\overset{\prime\prime}{,}}8\) in a triangle in which one direction also went to an extremely hard-to-see
\par\smallskip\noindent\rule{.24\textwidth}{.4pt}\par\smallskip
{\footnotesize\noindent\hspace*{1.5em}\textsuperscript{*}) He also has 3 heliotropes of the second design, made by Mr.~\textsc{Rumpf}.\par}
\clearpage
\noindent\null\hfill{\small\textbf{ON THE HANOVER TRIANGULATION.}}\hfill\makebox[0pt][r]{367}\par\medskip
\noindent non-heliotropic target.
I had wished to visit the station in the latter case once more and take the direction by heliotrope light, but could no longer get to it this year.
The large triangle Hohehagen-Brocken-Inselsberg is not included among the 26. The angle on Inselsberg was measured by \textsc{Gerling}, and the weather was very unfavourable for my measurements on Brocken, so that I could sight Inselsberg only 15 times (1823). Combined with the 15 sightings of 1821, which differ from them by less than \(1^{\prime\prime}\), they nevertheless give an excellent closure for this large triangle. --- In a country like Württemberg or Electoral Hesse, where there are so many high points, surveying is a pleasure and the largest triangles are easy to find.
It is the same in the southern part of Hanover.
But in the north, on the Lüneburg Heath, I have had unspeakable difficulties, and a reconnaissance undertaken last summer westward towards Bremen has so far yielded nothing to suggest even passably good triangles are possible.
High scaffolds could well overcome the difficulties, but I am apprehensive about the great cost in money and time, and even more about the loss of stability in the setup.
Farther south it would presumably be easier to get through the Osnabrück region to Bentheim, but most of the measurements would then have to extend across foreign territory.

\begin{center}\rule{.24\textwidth}{.4pt}\end{center}
